\section{Induction over the original composition}\label{cs:assembly}

We now explain why the two estimates just proved exhaust a strict
source, and why their exponents add to its own complete action.
The reduction is an equality of finite contractions. The original
ordered diagram is still needed for any later geometric interpretation.

\subsection{The reduced tree and its constants}

Associate consecutive products and cancel pairs of adjacent turns
using \(\Turn^2=1\) and \eqref{cs:adjoint}. The result is a tree of
product junctions, each of valence at least three. Every internal edge
carries a turn. An external edge carries either
\(h_i=\Turn(R^{\epsilon_i}L_{j_i})\), or its turn
\(p_i=R^{\epsilon_i}L_{j_i}\). Call the latter edge \emph{flagged}.
The reduced-junction theorem in \cite{Profiles} says that strictness
is equivalent to at most one flagged external edge at each junction.
All labels and full polynomial classes remain in this tree.

A junction with a flagged leaf sends toward its parent an operator
\begin{equation}\label{cs:L-module}
 F_i=D_{j_i}^{\epsilon_i}[E_i],\qquad E_i=\prod_b Q_b.
\end{equation}
An ordinary external leaf is the same operation with \(E_i=1\).
These are the type-\(L\) modules used at the boundaries of
Section~\ref{cs:probe}. Every maximal connected set of unflagged
junctions gives the other operation,
\begin{equation}\label{cs:A-module}
 Q=\Turn A(F_1,\ldots,F_n),\qquad n\ge2.
\end{equation}
Its boundary modules are all type \(L\); adjoining unflagged junctions
were already included in the same maximal component. Both operations
depend only on smaller subtrees.

Let \(f_Q(b)\) be the optimized pure action with output \(b\) held
fixed, and define its intrinsic rate
\[
 \Omega_Q=\max_{0\le b\le1/2}\{g(b)+f_Q(b)\}.
\]
The harmonic-profile and probe theorems of \cite{Profiles} give
\begin{align}
 f_{F_i}(b)&=g(b)+c_i,
 &c_i&=\sum_b\Omega_{Q_b},
 &\Omega_{F_i}&=c_i+2G_0,\label{cs:rigid-budget}\\
 \Omega_{\Turn A(F_1,\ldots,F_n)}
     &=\sum_i c_i+\Phi_A.&&&&\label{cs:component-budget}
\end{align}
The sum in \eqref{cs:rigid-budget} is zero for an empty input.
These equalities include the additive constants, not just derivatives
of profiles. The quantity in \eqref{cs:component-budget} is precisely
\eqref{cs:Omega}. Thus the positive comparison has not introduced
an extra exponential budget.

\begin{figure}[t]
\centering\begin{tikzpicture}[>=Latex,thick,font=\small,
 junction/.style={circle,draw,fill=blue!6,minimum size=6mm,inner sep=1pt},
 factor/.style={draw,rounded corners=2pt,inner sep=5pt}]
\node[junction] (a) at (0,1) {$\times$};
\node[junction] (b) at (0,-.5) {$\times$};
\node[factor] (f1) at (-1.5,1) {$F_1$};
\node[factor] (f2) at (-1,-1.8) {$F_2$};
\node[factor] (f3) at (1,-1.8) {$F_3$};
\draw[->] (f1)--(a);
\draw[->] (f2)--(b);
\draw[->] (f3)--(b);
\draw[->] (b)--node[right]{$\Turn$}(a);
\draw[->] (a)--node[right]{$\Turn$}(0,2.2);
\node[above] at (0,2.2) {$Q=\Turn A$};
\draw[gray!50] (2.1,-2.2)--(2.1,2.5);
\node[align=left,anchor=west] at (2.6,1.5)
 {$\displaystyle W_N=\sum_r w_{N,r}e_r$};
\node[align=left,anchor=west] at (2.6,.35)
 {$\displaystyle \normJ{Q_N}\le CN^A\phi(W_N A_N^{\mathrm{abs}})$};
\node[align=left,text width=6.4cm,anchor=west] at (2.6,-1.1)
 {$A_N^{\mathrm{abs}}$ uses the same products and turns, with the absolute
 values of the boundary coefficients.};
\end{tikzpicture}

\caption{An unflagged region with three boundary modules.
Here \(A=F_1\Turn(F_2F_3)\) and \(Q=\Turn A\). The internal
turn is part of the original expression. The positive scalar on the
right bounds its outgoing norm with the same shared indices.
The polynomial \(W_N\) belongs to this inequality only.}
\label{cs:component-figure}
\end{figure}

\subsection{Properties propagated from the leaves}

We prove three properties together. Every subtree output \(Q\) has
the norm bound \(e^{N(\Omega_Q+o(1))}\), a canonical angular strip
bound at that same rate with a strictly negative angular quadratic,
and a fixed norm loss whenever any one of its original colors has a
macroscopic deficit. In addition, a type-\(L\) module \(F_i\) has
the coefficient envelope \eqref{cs:cone}, and each original color
inside that module gives a strict coefficient loss on every fixed
compact interval in \((0,1/2)\).

For an elementary leaf, these assertions are
Lemma~\ref{cs:duality}, Proposition~\ref{cs:decorated-turn} and
Proposition~\ref{cs:regular-deficit}, with \(E=1\).
In particular its coefficient reference constant is zero and its
norm rate is \(2G_0\).

Suppose next that \(F_i\) is given by \eqref{cs:L-module}.
Multiply the canonical representatives of the \(Q_b\) without first
reducing that product. The norm bounds multiply by
submultiplicativity, their analytic strip bounds multiply as actual
polynomials, and their rates sum to \(c_i\) in
\eqref{cs:rigid-budget}. A deficit in a descendant gives a strict
norm loss in this product. Lemma~\ref{cs:duality} then preserves
that loss in every output coefficient. A deficit in the newly
inserted character gives the regular-index coefficient loss by
Proposition~\ref{cs:regular-deficit}. That proposition also supplies
the norm loss for the new color. Proposition~\ref{cs:decorated-turn}
gives the canonical strip estimate at rate \(c_i+2G_0=\Omega_{F_i}\).
This proves every claimed property for the new type-\(L\) module.

For \eqref{cs:A-module}, the input coefficient envelopes give
Theorem~\ref{cs:unflagged}, and Proposition~\ref{cs:probe-strip}
regenerates the angular estimate at \(\Omega_Q\). A deficit in any
original descendant lies inside one boundary module \(F_i\).
Its coefficient loss holds on a compact interval about the actual
selected index \(a_i^*\), even when that index is above \(1/4\).
Corollary~\ref{cs:loss-transfer} gives the required norm loss in
\(Q\). Equation \eqref{cs:component-budget} identifies its rate
with the intrinsic rate of this very subtree.

This is an induction on a finite tree. At each step the constants may
depend on the subtree, but not on its physical color tuple. Taking
the minimum of finitely many positive losses leaves a positive loss.
No lower estimate or nonzero value of an arbitrary subtree is assumed.

\subsection{The closing character, including an unflagged root}

The arbitrary-leaf rooting identity of \cite{Profiles} allows one
original leaf to be exposed without changing the contraction. If its
edge parity is \(\sigma\in\{0,1\}\), and \(A\) is the product of
the other branches at its junction, then
\begin{equation}\label{cs:any-root}
 Z_N=\phi((\Turn^\sigma h_0)A)
     =\phi(R^{\epsilon_0}L_{j_0}\,\Turn^{\sigma+1}A).
\end{equation}
There are at least two other branches because \(H\ge3\) and
junctions have valence at least three.

If \(\sigma=1\), the exposed leaf is the sole possible flag there.
The final input in \eqref{cs:any-root} is the unturned product of
the remaining branches. Its norm rate is their sum of intrinsic
rates, which the root-budget identity of \cite{Profiles} identifies
with \(F_0\). Their strip estimates multiply to give a positive
angular quadratic at that rate.

If \(\sigma=0\) and another leaf at this junction is flagged, the
final input is the type-\(L\) module supplied by that other
character. If no leaf there is flagged, it is the outgoing turn
of an unflagged component. Both cases have already been treated
by the induction, including their angular estimates. The same
root-budget identity says \(\Omega_{\Turn A}=F_0\). One can also
see this equality directly from the positive closing constraint:
\[
 F_0=\max_a\{g(a)+f_A(1-a)\}=\Omega_{\Turn A}.
\]

Now apply \eqref{cs:finite-cost} to bound the full contraction.
A deficit in a descendant gives a strict norm loss by the induction;
the closing character has only polynomial norm cost. A deficit in
the exposed color \(j_0\) is handled by
Theorem~\ref{cs:root-bound} on the inherited strip. This proves
\eqref{cs:whole} and \eqref{cs:color-gap}. It also treats a tree with no
flagged junction anywhere. In particular the three-leaf closure
uses exactly the \(n=2\) comparison in
Theorem~\ref{cs:unflagged}; it requires no additional source theorem.

\subsection{From a physical color loss to the core comparison}

For a physical high color \(j>N/2\), the exact circle formula gives
\begin{equation}\label{cs:physical-alias}
 h_j^\epsilon=C_j^\epsilon+(-1)^N h_{N-j}^\epsilon,
       \qquad h_0^\epsilon=0.
\end{equation}
Here \(C_j^\epsilon\) is the selected core; set it equal to zero
for \(j\le N/2\). Formula \eqref{cs:physical-alias} follows from
the pure complementary-color identity of \cite[Section 4]{Newton}
by the linear operation \(F\mapsto\Turn(R^\epsilon\Turn F)\).
It therefore retains the crossing and all repeated-root derivatives.

If some original color is at most \(N/2\), the all-core contraction
is zero, and \eqref{cs:color-gap} with \(\delta=1/2\) proves
\eqref{cs:alias-gap}. Otherwise replace each core in its
multilinear contraction by
\(h_j^\epsilon-(-1)^N h_{N-j}^\epsilon\).
Subtracting from the full source leaves finitely many terms with
at least one color \(N-j\le N/2\). Terms containing color zero
vanish. Every other term is an original physical source with a
macroscopic color deficit, so \eqref{cs:color-gap} applies. The number
of terms is at most \(2^H-1\), independent of \(N\). Reducing
the gap absorbs that finite sum and proves
\eqref{cs:alias-gap}, completing Theorem~\ref{cs:main}.

\subsection{The strict two-leaf closure}

\begin{corollary}[Two leaves with their original signs]\label{cs:two-leaf-source}
Let \(H=2\) and suppose the closure is strict. After the exact turn
and pairing identities it is
\[
 Z_N=\phi(h_{j_0}^{\epsilon_0}\Turn h_{j_1}^{\epsilon_1})
     =\phi(p_0\Turn p_1),\qquad
 p_i=R^{\epsilon_i}L_{j_i},\quad \epsilon_i\in\{0,1,-1\}.
\]
Its pure complete rate is \(F_0=2G_0\). The whole-source bound
\eqref{cs:whole}, every macroscopic physical-color gap
\eqref{cs:color-gap}, and the full-minus-core gap
\eqref{cs:alias-gap} hold with this value, at all sufficiently large
integer levels. No unflagged-component comparison is needed.
\end{corollary}
\begin{proof}
Strictness leaves one turn between the two elementary leaves, modulo
\(\Turn^2=1\) and the adjoint pairing identity. The displayed equality
then follows from \(h_j^\epsilon=\Turn p_j\).
For the pure complete circle, the exact coefficients in
\eqref{cs:split} give
\[
 h_N=\sum_{r=0}^M a_{N,r}e_r,\qquad
 a_{N,r}=\frac{N^2S_{2r}}{S_r^4},\qquad
 \phi(h_N\Turn h_N)=\sum_{r=0}^M a_{N,r}^2.
\]
Thus the normalization is the original coefficient-square sum, without
an extra trace factor. Uniform factorial rates and one index nearest
\(N/4\) give its all-level exponential rate \(2G_0\).

The final input to the character \(p_0\) is exactly
\(\Turn p_1=D_{j_1}^{\epsilon_1}[1]\). Proposition~\ref{cs:decorated-turn}
gives its norm and angular strip at rate \(2G_0\). A deficit in
\(j_1\) gives a norm gap by Proposition~\ref{cs:regular-deficit};
the other character then has only the polynomial cost
\eqref{cs:finite-cost}. A deficit in \(j_0\) gives a gap by
Theorem~\ref{cs:root-bound} on that same strip. The norm bound also
gives the full upper rate. Finally the exact physical alias expansion
from the preceding subsection converts these color gaps into the
full-minus-core gap. It preserves both retained crossing signs and
all double-root derivatives.
\end{proof}
